%%
%% This is file `poli-quickref.tex',
%% generated with the docstrip utility.
%%
%% The original source files were:
%%
%% ufrj-poli.dtx  (with options: `quickrefpoli')
%% 
%% This is a sample undergraduate project which illustrates the use of the
%% `ufrj' document class with the `ufrj-poli' unit style. It is generated from
%% ufrj-poli.dtx by docstrip.
%% 
%% Copyright (C) 2008-2026 Vicente Helano F. Batista, George O. Ainsworth Jr.,
%% Geraldo Xexéo and Eduardo Mangeli. Released under the GNU General Public
%% License version 3 (see the COPYING.txt file).
%% 
%% poli-quickref.tex -- CoppeTeX quick reference, Escola Politecnica, in English.
%%
%% GERADO de ufrj-poli.dtx. Nao edite este arquivo: edite o modulo
%% `quickrefpoli' do ufrj-poli.dtx e rode `pdflatex ufrj-poli.ins'.
\documentclass[a4paper,10pt]{article}
\usepackage[T1]{fontenc}
\usepackage[utf8]{inputenc}
\usepackage{lmodern}
\usepackage[margin=2cm]{geometry}
\usepackage{longtable}
\usepackage{booktabs}
\usepackage{array}
\usepackage{xcolor}
\usepackage[colorlinks=true,allbordercolors=white]{hyperref}

\newcommand\ufrjversion{v5.0}

\newcolumntype{C}{>{\ttfamily\raggedright\arraybackslash}p{0.17\textwidth}}
\newcolumntype{D}{>{\raggedright\arraybackslash}p{0.44\textwidth}}
\newcolumntype{T}{>{\raggedright\arraybackslash}p{0.30\textwidth}}

\setlength\parindent{0pt}

\begin{document}

\begin{center}
  {\Large\bfseries CoppeTeX quick reference --- Escola Politécnica}\\[2pt]
  {\normalsize Unit style \texttt{ufrj-poli}, \ufrjversion}\\[6pt]
  \begin{minipage}{0.92\textwidth}\small
  This sheet is the Escola Politécnica half. The other half --- every option,
  command and environment of the class --- is \texttt{ufrj-quickref.pdf}; the
  full manual of the style is \texttt{ufrj-poli.pdf}, in Portuguese.

  \medskip
  A \emph{Projeto de Graduação} is the class with \texttt{grad} plus this
  style:

  \medskip
  \verb|\documentclass[grad]{ufrj}|\\
  \verb|\usepackage{ufrj-poli}|

  \medskip
  The style brings the school's name and colour mark to the cover, beside
  UFRJ's, its Courses, the phrase ``Projeto de Graduação apresentado ao Curso
  de Engenharia X'' and the title each Course confers. Everything else is the
  class.
  \end{minipage}
\end{center}

\vspace{6pt}

\begin{longtable}{@{}CTD@{}}
\toprule
\normalfont\bfseries \textbackslash department & \normalfont\bfseries Course
& \normalfont\bfseries Title conferred \\
\midrule
\endfirsthead
\toprule
\normalfont\bfseries \textbackslash department & \normalfont\bfseries Course
& \normalfont\bfseries Title conferred \\
\midrule
\endhead
\bottomrule
\endfoot
AMBIENTAL      & Engenharia Ambiental & Engenheiro Ambiental \\
CIVIL          & Engenharia Civil & Engenheiro Civil \\
COMPUTACAO     & Engenharia de Computação e Informação & Engenheiro de Computação e Informação \\
CONTROLE       & Engenharia de Controle e Automação & Engenheiro de Controle e Automação \\
ELETRICA       & Engenharia Elétrica & Engenheiro Elétrico \\
ELETRONICA     & Engenharia Eletrônica e de Computação & Engenheiro Eletrônico e de Computação \\
MATERIAIS      & Engenharia de Materiais & Engenheiro de Materiais \\
MECANICA       & Engenharia Mecânica & Engenheiro Mecânico \\
METALURGICA    & Engenharia Metalúrgica & Engenheiro Metalúrgico \\
NAVAL          & Engenharia Naval e Oceânica & Engenheiro Naval \\
NUCLEAR        & Engenharia Nuclear & Engenheiro Nuclear \\
PETROLEO       & Engenharia de Petróleo & Engenheiro de Petróleo \\
PRODUCAO       & Engenharia de Produção & Engenheiro de Produção \\
NANOTECNOLOGIA & Nanotecnologia & the Course's own (it confers no engineer's title) \\
\midrule
\multicolumn{3}{@{}l@{}}{\normalfont\bfseries\small THE OLD DEPARTMENT ACRONYMS,
  from \texttt{poli.cls}, all still valid}\\[2pt]
ECI  & Engenharia de Computação e Informação & Engenheiro de Computação e Informação \\
DEL  & Engenharia Eletrônica e de Computação & Engenheiro Eletrônico e de Computação \\
DEE  & Engenharia Elétrica & Engenheiro Elétrico \\
DEM  & Engenharia Mecânica & Engenheiro Mecânico \\
DMM  & Engenharia Metalúrgica e de Materiais & Engenheiro Metalúrgico e de Materiais \\
DNC  & Engenharia Nuclear & Engenheiro Nuclear \\
DENO & Engenharia Naval e Oceânica & Engenheiro Naval \\
DEI  & Engenharia de Produção & Engenheiro de Produção \\
DCC, DES, DEG, DET, DHIMA & Engenharia Civil, or the department's own name with the option \texttt{civilpordepartamento} & Engenheiro Civil \\
\end{longtable}

\vspace{4pt}
\begin{center}
\begin{minipage}{0.92\textwidth}\small
\textbf{Options of the style.}
\verb|civilporcurso| (the default) makes the five Civil acronyms read
``Engenharia Civil'', which is the Course; \verb|civilpordepartamento| makes
each read its department's name, as \texttt{poli.cls} did.
\verb|times| and \verb|arial| set the family of the letter, which Resolução
05/2012 leaves to the author --- the SIZE stays 12 in all three, which is what
the Manual fixes (2.2b); without either, the class's family is used.

\medskip
\textbf{The school's own rule.} Resolução 05/2012 is older than the Manual and
contradicts it in thirteen points --- pre-textual folios in roman numerals, the
sumário before the lists, the caption below the illustration, references based
on NB-66, an ABNT standard replaced in 1989. Where they disagree the Manual
wins, and the style implements only what the Resolução \emph{chooses} among
what the Manual leaves open. The check, item by item, is in
\texttt{NAO-CONFORMIDADES-POLI.md}, and the minimal norm proposed in its place
is \texttt{PROPOSTA-DE-RESOLUCAO.pdf}. \textbf{The bibliography follows the
UFRJ rule strictly}; any difference presented in the Resolução is wrong.

\medskip
\textbf{A project begun with the old class} compiles as it is: \texttt{poli.cls}
loads \texttt{ufrj} with \texttt{grad} and this style, accepts the dead options
(\texttt{pdftex}, \texttt{dvips}, \texttt{dvipdfm}, \texttt{draft}) and the old
signatures of \verb|\advisor|, \verb|\coadvisor| and \verb|\examiner|, warning
once that the board comes out without the institution that 3.1.2.1.3(e) asks
for. It relaxes no rule of the class --- the order of the pre-textual elements
is an error there too.

\medskip
Worked examples: \texttt{poli-min-exemplo.tex}, only what is mandatory, and
\texttt{poli-max-exemplo.tex}, everything the class composes.\quad
Problems: \url{https://github.com/COPPE-UFRJ/CoppeTeX/issues}
\end{minipage}
\end{center}

\end{document}
%% 
%%
%% End of file `poli-quickref.tex'.
